% Fuente: Auxiliar 4, Problema 1.
{\color{MidnightBlue}
\begin{enumerate}

\item
Por definición, $x \in P$ es un vértice de $P$ si existe un vector $c \in \mathbb{R}^n$ tal que
    \[
        c^\top x < c^\top u, \qquad \forall u \in P, \; u \neq x.
    \]

    Supongamos, por contradicción, que $x$ no es un punto extremo de $P$. Entonces existen
    $y,z \in P$, con $y \neq z$, y un escalar $\lambda \in (0,1)$ tales que
    \[
        x = \lambda y + (1-\lambda)z.
    \]
    Además, como $x$ no es extremo, podemos tomar $y \neq x$ y $z \neq x$. Luego, por la
    definición de vértice,
    \[
        c^\top x < c^\top y,
        \qquad
        c^\top x < c^\top z.
    \]
    Multiplicando por $\lambda$ y $(1-\lambda)$, respectivamente, se obtiene
    \[
        \lambda c^\top x < \lambda c^\top y,
        \qquad
        (1-\lambda)c^\top x < (1-\lambda)c^\top z.
    \]
    Sumando ambas desigualdades,
    \[
        c^\top x
        <
        \lambda c^\top y + (1-\lambda)c^\top z.
    \]
    Sin embargo, usando que $x = \lambda y + (1-\lambda)z$,
    \[
        \lambda c^\top y + (1-\lambda)c^\top z
        =
        c^\top\left(\lambda y + (1-\lambda)z\right)
        =
        c^\top x.
    \]
    Por lo tanto,
    \[
        c^\top x < c^\top x,
    \]
    lo cual es una contradicción. Luego, $x$ debe ser un punto extremo de $P$.

\item
Como $A \in \mathbb{R}^{n \times n}$ tiene filas linealmente independientes, entonces $A$ es no singular.
    Por lo tanto, el sistema
    \[
        Ax=b
    \]
    tiene una única solución, dada por
    \[
        x^\star = A^{-1}b.
    \]

    La factibilidad del problema depende de si esta solución satisface la restricción de no negatividad.
    Por lo tanto:

    \begin{itemize}
        \item Si
        \[
            A^{-1}b \geq 0,
        \]
        entonces el problema es factible. Además, como existe un único punto factible,
        \[
            x^\star = A^{-1}b,
        \]
        este punto es necesariamente la única solución óptima del problema, para cualquier vector de costos $c$.
        El valor óptimo es
        \[
            z^\star = c^\top A^{-1}b.
        \]

        \item Si existe alguna componente negativa en $A^{-1}b$, es decir,
        \[
            (A^{-1}b)_i < 0
        \]
        para algún $i$, entonces el problema es infactible, porque la única solución de $Ax=b$
        no satisface $x \geq 0$. En este caso, el problema no tiene solución óptima.
    \end{itemize}

    En conclusión, el problema tiene solución óptima si y solo si
    \[
        A^{-1}b \geq 0.
    \]
    En tal caso, existe una única solución óptima:
    \[
        x^\star = A^{-1}b.
    \]

\item
Un ejemplo es el siguiente problema:
    \[
        \begin{aligned}
            \min_{x_1,x_2} \quad & x_1+x_2 \\
            \text{s.a.} \quad & x_1 + x_2 \geq 1, \\
            & x_1,x_2 \in \mathbb{R}.
        \end{aligned}
    \]
    Su región factible es
    \[
        P = \{(x_1,x_2) \in \mathbb{R}^2 \mid x_1+x_2 \geq 1\}.
    \]
    Este conjunto es no vacío, ya que, por ejemplo,
    \[
        (0,1) \in P.
    \]

    Ahora probemos que $P$ no contiene puntos extremos. Basta con demostrar que contiene una linea. Sea $x=(x_1,x_2) \in P$ y considere
    la dirección
    \[
        d = (1,-1).
    \]
    Para cualquier $\varepsilon > 0$, se tiene
    \[
        x+\varepsilon d = (x_1+\varepsilon, x_2-\varepsilon),
    \]
    y
    \[
        x-\varepsilon d = (x_1-\varepsilon, x_2+\varepsilon).
    \]
    En ambos casos, la suma de las componentes se mantiene igual:
    \[
        (x_1+\varepsilon)+(x_2-\varepsilon) = x_1+x_2,
    \]
    y
    \[
        (x_1-\varepsilon)+(x_2+\varepsilon) = x_1+x_2.
    \]
    Como $x \in P$, se cumple que $x_1+x_2 \geq 1$. Por lo tanto,
    \[
        x+\varepsilon d \in P,
        \qquad
        x-\varepsilon d \in P.
    \]

    En consecuencia, la linea
    \[
    \{x+td\:|\:d=(1,-1) \: | \:y \in \mathbb{R}\}\:\in \:P
    \]
    esta contenida en $P$. Por lo tanto $P$ es un conjunto no vacío que no contiene puntos extremos.

\end{enumerate}
}

